% Section: Observables
\section{Observables}
\label{sec:observables}

This section defines the observables computed by the analysis tasks of \CAP{} and lists, for each task, the histograms
it produces. Three conventions apply to all of them.

\paragraph{Acceptance.}
The kinematic range and the species of the particles entering an observable are those of the particle filters of the
task (Section~\ref{sec:filters}); single-particle and pair histograms are incremented only for the particles accepted
by their filters. The binning of each histogram is set by the histogram settings of the task (keys
\code{<task>:HISTOGRAM:...}), and its range should cover the range accepted by the filters.

\paragraph{Normalisation.}
At the end of the event loop, the histograms of the single-particle and pair analyses are divided by the number of
events accepted by their event filter: they hold yields \emph{per event}. Bin contents are integrals over the bins, not
densities: a histogram of $\pT$ holds $\avg{\Delta N}$, the mean number of particles per event in each $\pT$ bin, and
the density $dN/d\pT$ is obtained by dividing by the bin width. Cross sections follow from the yields per event by
multiplication with the cross section of the events selected, $d\sigma = \sigma_{\mathrm{ev}}\, dN$. Distributions of
event-wise quantities (multiplicity, energy, spherocity) and profiles are not scaled.

\paragraph{Base and derived histograms.}
Each table distinguishes the \emph{base} histograms, filled event by event by the analyzer, from the \emph{derived}
histograms, computed from them by the derived calculator in a later job (Section~\ref{sec:phases}). The last column
gives the option of the histogram settings that turns a histogram on: \code{eta\_fill}, \code{rapidity\_fill},
\code{p2\_fill}, \code{FillPtvsY}, \code{FillPid}; histograms without option are always produced. The histograms of a
combination of filters are in its folder of the output file (e.g. \code{ALL\_PiP} for singles, \code{ALL\_PiP\_PiM}
for pairs).

% ---------------------------------------------------------------------------------------------------------------
\subsection{Single-particle spectra}
\label{sec:single}

The task \code{ParticleSingleAnalyzer} measures the single-particle densities of each species selected by a particle
filter,
\begin{equation}
  \rho_1(\pT,\eta,\varphi) = \frac{d^3N}{\pT\,d\pT\,d\eta\,d\varphi},
\end{equation}
and its projections, and the event-by-event distributions of the multiplicity $N$ and total energy $E$ of the accepted
particles. The invariant yield, and the invariant cross section $E\,d^3\sigma/dp^3 = \sigma_{\mathrm{ev}}\,E\,d^3N/dp^3$,
\begin{equation}
  E\frac{d^3N}{dp^3} = \frac{1}{2\pi\pT}\frac{d^2N}{d\pT\,dy},
  \label{eq:invariant}
\end{equation}
is obtained from \code{n1\_ptXS}, filled with weight $1/\pT$, as $\code{n1\_ptXS}/(2\pi\,\Delta\pT\,\Delta y)$, with
$\Delta\pT$ the width of the bin and $\Delta y$ the rapidity range accepted by the particle filter; the double
differential $d^2N/d\pT\,dy$ is available from \code{n1\_ptY}. The azimuthal and (pseudo)rapidity densities are
\begin{equation}
  \rho_1(\eta,\varphi) = \frac{d^2N}{d\eta\,d\varphi}, \;\;
  \frac{dN}{d\eta} = \int\! \rho_1\, d\varphi, \;\;
  \frac{dN}{d\varphi} = \int\! \rho_1\, d\eta ,
\end{equation}
and likewise in rapidity. When \code{p2\_fill} is set, the same histograms are also filled with weight $\pT$, which
gives the inclusive mean transverse momentum in each bin, $\avg{\pT}(\eta,\varphi) = \sum\pT/N$. Table~\ref{tab:single}
lists the histograms.

\begin{table*}[t]
  \centering
  \caption{Histograms of a single-particle folder (\code{ParticleSingleHistos}), per event.}
  \label{tab:single}
  \small
  \begin{tabular}{@{}lll@{}}
    \toprule
    Name & Content & Option \\
    \midrule
    \multicolumn{3}{@{}l}{\emph{Base}} \\
    \code{n1}           & distribution of $N$             & \\
    \code{n1\_eTotal}   & distribution of $E=\sum E_i$    & \\
    \code{n1\_pt}       & $dN/d\pT$                        & \\
    \code{n1\_ptXS}     & $(1/\pT)\,dN/d\pT$               & \\
    \code{n1\_phiEta}   & $d^2N/d\eta\,d\varphi$           & \code{eta\_fill} \\
    \code{spt\_phiEta}  & $\sum\pT$ vs $(\eta,\varphi)$    & \code{eta\_fill}, \code{p2\_fill} \\
    \code{n1\_phiY}     & $d^2N/dy\,d\varphi$              & \code{rapidity\_fill} \\
    \code{spt\_phiY}    & $\sum\pT$ vs $(y,\varphi)$       & \code{rapidity\_fill}, \code{p2\_fill} \\
    \code{n1\_ptY}      & $d^2N/dy\,d\pT$                  & \code{FillPtvsY} \\
    \code{n1\_indexId}  & PDG codes of the particles       & \code{FillPid} \\
    \midrule
    \multicolumn{3}{@{}l}{\emph{Derived}} \\
    \code{n1\_phi}      & $dN/d\varphi$                    & \\
    \code{n1\_eta}      & $dN/d\eta$                       & \code{eta\_fill} \\
    \code{n1\_y}        & $dN/dy$                          & \code{rapidity\_fill} \\
    \code{sum\_pt\_}$X$ & $\sum\pT$ vs $X=\varphi,\eta,y$  & \code{p2\_fill} \\
    \code{avg\_pt\_}$X$ & $\avg{\pT}$ vs $X=\varphi,\eta,y,(\eta,\varphi),(y,\varphi)$ & \code{p2\_fill} \\
    \bottomrule
  \end{tabular}
\end{table*}

% ---------------------------------------------------------------------------------------------------------------
\subsection{Pair densities and cumulants}
\label{sec:pairs}

The task \code{ParticlePairAnalyzer} measures, for each ordered pair of particle filters $(\alpha,\beta)$, the pair
density $\rho_2^{\alpha\beta}(x_1,x_2)$ as a function of the kinematic variables of the two particles, $x=\pT$,
$\varphi$, $\eta$ or $y$, and of their differences $(\Delta\eta,\Delta\varphi)$ and $(\Delta y,\Delta\varphi)$, together with
the single-particle densities of the two species (Section~\ref{sec:single}). When $\alpha=\beta$, a particle is not
paired with itself. The derived histograms are the product of the single-particle densities $\rho_1^\alpha\otimes\rho_1^\beta$,
the cumulant, and the normalised cumulant~\cite{Pruneau:2002yf,Pruneau:2017book},
\begin{align}
  C_2^{\alpha\beta}(x_1,x_2) & = \rho_2^{\alpha\beta}(x_1,x_2) - \rho_1^{\alpha}(x_1)\,\rho_1^{\beta}(x_2), \label{eq:C2} \\
  R_2^{\alpha\beta}(x_1,x_2) & = \frac{C_2^{\alpha\beta}(x_1,x_2)}{\rho_1^{\alpha}(x_1)\,\rho_1^{\beta}(x_2)}, \label{eq:R2}
\end{align}
and the cumulant per associated particle,
\begin{equation}
  A_2^{\alpha\beta}(x_1,x_2) = \frac{C_2^{\alpha\beta}(x_1,x_2)}{\avg{N_\beta}},
  \label{eq:A2}
\end{equation}
with $\avg{N_\beta}$ the mean number of particles $\beta$ per event in the acceptance. For the differences of
pseudorapidity and azimuth, $\rho_1^{\alpha}\otimes \rho_1^{\beta}(\Delta \eta, \Delta \varphi)$ is the ``tensor
product'' of the single-particle densities $\rho_1^{\alpha}(\eta_1,\varphi_1)$ and $\rho_1^{\beta}(\eta_2,\varphi_2)$
averaged at $\Delta \eta= \eta_1-\eta_2$ and $\Delta \varphi= \varphi_1-\varphi_2$ according to the technique described
in Ref.~\cite{Adam:2017ucq}, and
\begin{equation}
  R_2^{\alpha\beta}(\Delta \eta, \Delta \varphi) = \frac{C_2^{\alpha\beta}(\Delta \eta, \Delta \varphi)}{\rho_1^{\alpha}\otimes \rho_1^{\beta}(\Delta \eta, \Delta \varphi)},
\end{equation}
with $C_2^{\alpha\beta}(\Delta \eta, \Delta \varphi) = \rho_2^{\alpha\beta}(\Delta \eta, \Delta \varphi) - \rho_1^{\alpha}\otimes \rho_1^{\beta}(\Delta \eta, \Delta \varphi)$,
and likewise for $(\Delta y,\Delta\varphi)$. $R_2$ is thereby averaged over the mean pseudorapidity
$\bar\eta=(\eta_1+\eta_2)/2$ of the pairs, and corrected for the acceptance, since the pair density and the product of
the single densities are reduced the same way. $D_2(\Delta y,\Delta\varphi)$ is $A_2$ divided, in each $\Delta y$ bin,
by the fraction $\Omega(\Delta y)$ of the acceptance in $\bar y$ available at that $\Delta y$ (for a uniform acceptance
$|y|<Y$, $\Omega = 1-|\Delta y|/2Y$), that is, $A_2$ averaged over $\bar y$. The $\Delta\varphi$ axis covers
$[0,2\pi)$; the versions \code{\_shft} of the histograms are shifted to $[-\pi/2,3\pi/2)$, where the near side
($\Delta\varphi\approx 0$) and the away side ($\Delta\varphi\approx\pi$) appear whole. Table~\ref{tab:pairs} lists the
histograms.
\todo{\code{C2O\_DetaDphi} ($D_2$ vs $\Delta\eta$) is created but not computed today; only $D_2(\Delta y,\Delta\varphi)$ is.}

\begin{table*}[t]
  \centering
  \caption{Histograms of a pair folder (\code{ParticlePairHistos}), per event, for the variables
  $X$ = \code{ptpt} $(p_{\mathrm{T},1},p_{\mathrm{T},2})$ and \code{phiphi} $(\varphi_1,\varphi_2)$, always;
  \code{etaeta} $(\eta_1,\eta_2)$ and \code{DetaDphi} $(\Delta\eta,\Delta\varphi)$ with \code{eta\_fill};
  \code{yY} $(y_1,y_2)$ and \code{DyDphi} $(\Delta y,\Delta\varphi)$ with \code{rapidity\_fill}.
  The column $X$ lists the variables for which a histogram is made.}
  \label{tab:pairs}
  \small
  \begin{tabular}{@{}llll@{}}
    \toprule
    Name & Content & $X$ & Option \\
    \midrule
    \multicolumn{4}{@{}l}{\emph{Base}} \\
    \code{n2}             & distribution of the number of pairs                     & ---         & \\
    \code{n2\_}$X$        & $\rho_2$, pair yield                                     & all         & \\
    \code{ptpt\_}$X$      & $\rho_2$ weighted by $p_{\mathrm{T},1}p_{\mathrm{T},2}$, $S_{pp}$ & all but \code{ptpt} & \code{p2\_fill} \\
    \code{ptn\_}$X$, \code{npt\_}$X$ & $\rho_2$ weighted by $p_{\mathrm{T},1}$ ($S_{pn}$) and by $p_{\mathrm{T},2}$ ($S_{np}$) & all but \code{ptpt} & \code{p2\_fill} \\
    \midrule
    \multicolumn{4}{@{}l}{\emph{Derived}} \\
    \code{n1n1\_}$X$      & $\rho_1^\alpha\otimes\rho_1^\beta$                        & all         & \\
    \code{C2\_}$X$        & $C_2$, Eq.~(\ref{eq:C2})                                  & all         & \\
    \code{R2\_}$X$        & $R_2$, Eq.~(\ref{eq:R2})                                  & all         & \\
    \code{C2N\_}$X$       & $A_2$, Eq.~(\ref{eq:A2})                                  & all         & \\
    \code{C2O\_}$X$       & $D_2$, $A_2$ averaged over $\bar y$                       & \code{DetaDphi}, \code{DyDphi} & \\
    \code{DptDpt\_}$X$    & $\avg{\Delta p_{\mathrm{T},1}\Delta p_{\mathrm{T},2}}$, Eq.~(\ref{eq:DptDpt}) & all but \code{ptpt} & \code{p2\_fill} \\
    \code{P2\_}$X$        & $P_2$, Eq.~(\ref{eq:P2})                                  & all but \code{ptpt} & \code{p2\_fill} \\
    \code{pt1pt1\_}$X$    & $\avg{p_{\mathrm{T},1}}\avg{p_{\mathrm{T},2}}$             & all but \code{ptpt} & \code{p2\_fill} \\
    \code{G2\_}$X$        & $G_2$, Eq.~(\ref{eq:G2})                                  & all but \code{ptpt} & \code{p2\_fill} \\
    $Q$\code{\_}$X$\code{\_shft} & $Q$ = \code{C2}, \code{C2N}, \code{C2O}, \code{R2}, \code{DptDpt}, \code{P2}, \code{G2}, $\Delta\varphi$ shifted & \code{DetaDphi}, \code{DyDphi} & \\
    \bottomrule
  \end{tabular}
\end{table*}

% ---------------------------------------------------------------------------------------------------------------
\subsection{Transverse-momentum correlators}
\label{sec:ptcorr}

With the option \code{p2\_fill}, the pair analysis also measures the correlations of the transverse momenta of the
particles of a pair. With $S_{pp}$, $S_{pn}$ and $S_{np}$ the pair densities weighted by $p_{\mathrm{T},1}p_{\mathrm{T},2}$,
$p_{\mathrm{T},1}$ and $p_{\mathrm{T},2}$, and $\avg{\pT}_\alpha$ the inclusive mean transverse momentum of species
$\alpha$ in the acceptance, the transverse-momentum covariance of the pairs in a bin is
\begin{align}
  \avg{\Delta p_{\mathrm{T},1}\Delta p_{\mathrm{T},2}} = {} & \frac{1}{\rho_2}\Big( S_{pp} - \avg{\pT}_\beta S_{pn} \nonumber \\
   & - \avg{\pT}_\alpha S_{np} + \avg{\pT}_\alpha\avg{\pT}_\beta\,\rho_2 \Big),
  \label{eq:DptDpt}
\end{align}
with $\Delta\pT = \pT-\avg{\pT}$, and its dimensionless version is~\cite{Sharma:2008qr}
\begin{equation}
  P_2 = \frac{\avg{\Delta p_{\mathrm{T},1}\Delta p_{\mathrm{T},2}}}{\avg{\pT}_\alpha\avg{\pT}_\beta}.
  \label{eq:P2}
\end{equation}
The correlator $G_2$~\cite{Gavin:2006xd} compares the $\pT$-weighted pair density with the product of the
$\pT$-weighted single densities,
\begin{equation}
  G_2 = \frac{S_{pp} - S^{\alpha}_{p}\otimes S^{\beta}_{p}}{\rho_1^\alpha\otimes\rho_1^\beta},
  \label{eq:G2}
\end{equation}
with $S_p$ the single-particle densities weighted by $\pT$ (\code{spt\_phiEta}, \code{spt\_phiY}).

The task \code{PtPtAnalyzer} is meant for the event-wise version of these correlators: the moments
$Q_n=\sum_i p_{\mathrm{T},i}^n$ of the accepted particles of each event and the factorial moments of their multiplicity,
accumulated as profiles versus an event variable, from which the cumulants of $\pT$ of orders 2 to 4 are derived.
\todo{\code{PtPtAnalyzer} applies the filters but does not fill its histograms yet: describe it once it does, or leave
it out of this version.}

% ---------------------------------------------------------------------------------------------------------------
\subsection{Balance functions}
\label{sec:bf}

Balance functions~\cite{Bass:2000az} measure the density of particles $\bar\beta$ associated with a particle $\alpha$,
minus that associated with its antiparticle $\bar\alpha$,
\begin{equation}
  B^{\alpha\bar\beta}(x_1,x_2) = \frac{\rho_2^{\alpha\bar\beta}(x_1,x_2)}{\rho_1^{\alpha}(x_1)}
                                - \frac{\rho_2^{\bar\alpha\bar\beta}(x_1,x_2)}{\rho_1^{\bar\alpha}(x_1)} ,
  \label{eq:BFdef}
\end{equation}
so that its integral over $x_2$ is the number of particles $\bar\beta$ that balance, in the acceptance, the
conserved charge carried by a particle $\alpha$. Since $\rho_2^{\alpha\beta}/\rho_1^\alpha = \rho_1^\beta(1+R_2^{\alpha\beta})$ by Eq.~(\ref{eq:R2}), it is computed from the normalised cumulants of
the pairs and the density of the associated species,
\begin{equation}
  B^{\alpha\bar\beta} = \rho_1^{\bar\beta}(x_2)\left[R_2^{\alpha\bar\beta}(x_1,x_2) - R_2^{\bar\alpha\bar\beta}(x_1,x_2)\right],
\end{equation}
and, for the differences of pseudorapidity (or rapidity) and azimuth, as defined in Ref.~\cite{Pruneau:2023zhl},
\begin{widetext}
\begin{equation}
  \label{eq:BF}
  B^{\alpha \bar{\beta}}(\Delta \eta, \Delta \varphi) = \bar\rho_1^{\bar{\beta}}  \left[ R_2^{\alpha\bar{\beta}}(\Delta \eta, \Delta \varphi) - R_2^{\bar{\alpha}\bar{\beta}}(\Delta \eta, \Delta \varphi) \right],
\end{equation}
\end{widetext}
with $\bar\rho_1^{\bar\beta}$ the mean density of $\bar\beta$ over the acceptance. The labels $\alpha, \beta$ denote the
particle species of interest, whereas $\bar\alpha$ and $\bar\beta$ denote their respective antiparticles. The balance
function $B^{\bar\alpha\beta}$ is obtained by exchanging particles and antiparticles, and the combination
$B^{s} = (B^{\alpha\bar\beta}+B^{\bar\alpha\beta})/2$ is reported as well. The charge-independent and
charge-dependent combinations of the normalised cumulants are
\begin{align}
  R_2^{\mathrm{CI}} & = \tfrac14\left(R_2^{\alpha\beta}+R_2^{\alpha\bar\beta}+R_2^{\bar\alpha\beta}+R_2^{\bar\alpha\bar\beta}\right), \\
  R_2^{\mathrm{CD}} & = \tfrac12\left(R_2^{\alpha\bar\beta}+R_2^{\bar\alpha\beta}-R_2^{\alpha\beta}-R_2^{\bar\alpha\bar\beta}\right).
\end{align}

The task \code{ParticlePairBFCalculator} computes them in the third phase of a pair analysis, from the folders of a
\code{ParticlePairAnalyzer}, for every event filter and ordered pair of species $(\alpha,\beta)$. A species is a pair of
particle filters, a particle and its antiparticle, named by the key \code{SPECIES} (e.g.\
\code{PiP/PiM, KP/KM, PP/PM}); by default, the first half of the particle filters are the particles and the second half
their antiparticles. The results are written in a folder per event filter and pair of species, named after the
particles with the suffix \code{\_BF} (e.g.\ \code{ALL\_PiP\_KP\_BF}); Table~\ref{tab:bf} lists its histograms. As for
the other histograms, bin contents are integrals over the bins: the sum of $B^{\alpha\bar\beta}$ over the bins of $x_2$
is the number of balancing partners.

\begin{table*}[t]
  \centering
  \caption{Histograms of a balance-function folder (\code{ParticlePairBFHistos}), for the variables $X$ of
  Table~\ref{tab:pairs} filled by the pair analysis; $(\Delta\eta,\Delta\varphi)$ and $(\Delta y,\Delta\varphi)$ have
  \code{\_shft} versions.}
  \label{tab:bf}
  \small
  \begin{tabular}{@{}ll@{}}
    \toprule
    Name & Content \\
    \midrule
    \code{R2CI\_}$X$    & $R_2^{\mathrm{CI}}$ \\
    \code{R2CD\_}$X$    & $R_2^{\mathrm{CD}}$ \\
    \code{B\_a\_bB\_}$X$ & $B^{\alpha\bar\beta}$, Eq.~(\ref{eq:BF}) \\
    \code{B\_aB\_b\_}$X$ & $B^{\bar\alpha\beta}$ \\
    \code{B\_}$X$       & $B^{s}$ \\
    \bottomrule
  \end{tabular}
\end{table*}

% ---------------------------------------------------------------------------------------------------------------
\subsection{Three-dimensional pair analysis}
\label{sec:pair3d}

The tasks \code{ParticlePair3DAnalyzer}, \code{ParticlePair3DDerivedCalculator} and \code{ParticlePair3DBFCalculator} are
called exactly as their counterparts of the pair analysis; only their histograms differ. They measure the pair
densities as functions of the relative momentum of the pairs, the Lorentz-invariant
$Q_{\mathrm{inv}} = \sqrt{-(p_1-p_2)^2}$ and the three components $(q_{\mathrm{side}}, q_{\mathrm{out}}, q_{\mathrm{long}})$
of $\mathbf{p}_1-\mathbf{p}_2$ in the longitudinally comoving frame of the pair, along the beam (long), along the
transverse momentum of the pair (out), and perpendicular to both (side). The product of the single-particle densities
in these variables has no closed form: it is computed by Monte Carlo, drawing pairs of particles from the measured
densities $d^2N/dy\,d\pT$ of the two species (\code{n1\_ptY}, option \code{FillPtvsY}) with uniform azimuths
(key \code{nMcSimEvents}). $C_2$ and $R_2$ follow from Eqs.~(\ref{eq:C2}) and~(\ref{eq:R2}), and the associated yields
are $A^{\alpha|\beta} = C_2/\avg{N_\beta}$ and $A^{\beta|\alpha}=C_2/\avg{N_\alpha}$. The balance functions are
computed per trigger particle, from Eq.~(\ref{eq:BFdef}) integrated over the variables of the trigger,
\begin{equation}
  B^{\alpha\bar\beta}(q) = \frac{\rho_2^{\alpha\bar\beta}(q)}{\avg{N_\alpha}} - \frac{\rho_2^{\bar\alpha\bar\beta}(q)}{\avg{N_{\bar\alpha}}},
\end{equation}
with $q = Q_{\mathrm{inv}}$, $(q_{\mathrm{side}},q_{\mathrm{out}},q_{\mathrm{long}})$ or one of its components.
Table~\ref{tab:pair3d} lists the histograms.
\todo{references for the Bertsch--Pratt variables and for balance functions in relative momentum.}

\begin{table*}[t]
  \centering
  \caption{Histograms of a 3D pair folder (\code{ParticlePair3DHistos}) and of its balance-function folder
  (\code{ParticlePair3DBFHistos}), for $q$ = \code{Qinv}, \code{DeltaP} $(q_{\mathrm{side}},q_{\mathrm{out}},q_{\mathrm{long}})$,
  and the projections \code{DeltaPs}, \code{DeltaPo}, \code{DeltaPl}.}
  \label{tab:pair3d}
  \small
  \begin{tabular}{@{}lll@{}}
    \toprule
    Name & Content & $q$ \\
    \midrule
    \multicolumn{3}{@{}l}{\emph{Base}} \\
    \code{n2}            & distribution of the number of pairs   & --- \\
    \code{n2\_}$q$        & $\rho_2$                              & \code{Qinv}, \code{DeltaP} \\
    \midrule
    \multicolumn{3}{@{}l}{\emph{Derived}} \\
    \code{n2\_}$q$        & $\rho_2$, projections                  & \code{DeltaPs,o,l} \\
    \code{n1n1\_}$q$      & $\rho_1\otimes\rho_1$ (Monte Carlo)    & all \\
    \code{c2\_}$q$, \code{r2\_}$q$ & $C_2$, $R_2$                  & all \\
    \code{A12\_}$q$, \code{A21\_}$q$ & $A^{\alpha|\beta}$, $A^{\beta|\alpha}$ & all but \code{DeltaP} \\
    \code{n1\_}$i$\code{\_pt}, \code{n1\_}$i$\code{\_y}, \dots & densities of the Monte Carlo pairs & --- \\
    \midrule
    \multicolumn{3}{@{}l}{\emph{Balance functions}} \\
    \code{R2CI\_}$q$, \code{R2CD\_}$q$ & $R_2^{\mathrm{CI}}$, $R_2^{\mathrm{CD}}$ & all \\
    \code{B\_a\_bB\_}$q$, \code{B\_aB\_b\_}$q$, \code{B\_}$q$ & $B^{\alpha\bar\beta}$, $B^{\bar\alpha\beta}$, $B^s$ & all \\
    \bottomrule
  \end{tabular}
\end{table*}

% ---------------------------------------------------------------------------------------------------------------
\subsection{Fluctuations: $\nu_{\rm dyn}$ and higher moments}
\label{sec:nudyn}

The task \code{NuDynAnalyzer} measures the fluctuations of the multiplicities $N_1$ and $N_2$ of two particle
selections, as functions of an event variable $m$ (the multiplicity or another global quantity of the event, see
Section~\ref{sec:global}), from the factorial moments
\begin{equation}
  f_1^{i} = \avg{N_i}, \;\;
  f_2^{ii} = \avg{N_i(N_i-1)}, \;\;
  f_2^{12} = \avg{N_1N_2},
\end{equation}
accumulated as profiles versus $m$. The derived calculator computes the factorial cumulants and their normalised
versions,
\begin{equation}
  F_2^{ij} = f_2^{ij} - f_1^{i}f_1^{j}, \qquad R_2^{ij} = \frac{F_2^{ij}}{f_1^{i}f_1^{j}},
\end{equation}
and the dynamical fluctuations of the relative multiplicities~\cite{Pruneau:2002yf},
\begin{equation}
  \nudyn = R_2^{11} + R_2^{22} - 2R_2^{12},
\end{equation}
which vanishes for independently produced particles. The task has a folder per event filter and ordered pair of
particle filters $(i,j)$; the folders with $j<i$ are filled. Table~\ref{tab:nudyn} lists the histograms.

\begin{table*}[t]
  \centering
  \caption{Histograms of a fluctuation folder (\code{NuDynHistos}), as functions of the event variable $m$.}
  \label{tab:nudyn}
  \small
  \begin{tabular}{@{}ll@{}}
    \toprule
    Name & Content \\
    \midrule
    \multicolumn{2}{@{}l}{\emph{Base}} \\
    \code{mult}                  & distribution of $m$ \\
    \code{f1\_1}, \code{f1\_2}     & $f_1^1$, $f_1^2$ (profiles) \\
    \code{f2\_11}, \code{f2\_12}, \code{f2\_22} & $f_2^{11}$, $f_2^{12}$, $f_2^{22}$ (profiles) \\
    \midrule
    \multicolumn{2}{@{}l}{\emph{Derived}} \\
    \code{F2\_11}, \code{F2\_12}, \code{F2\_22} & $F_2^{ij}$ \\
    \code{R2\_11}, \code{R2\_12}, \code{R2\_22} & $R_2^{ij}$ \\
    \code{Nu\_dyn}               & $\nudyn$ \\
    \bottomrule
  \end{tabular}
\end{table*}
\todo{the code computes moments up to second order; higher moments (orders 3 and 4) are announced in the documentation
of the class but not implemented.}

% ---------------------------------------------------------------------------------------------------------------
\subsection{Anisotropic flow}
\label{sec:flow}

The task \code{ParticleFlowAnalyzer} measures the Fourier coefficients $v_n$ of the azimuthal distributions of the
particles, $dN/d\varphi \propto 1 + 2\sum_n v_n\cos n(\varphi-\Psi_n)$, by multi-particle cumulants, for the harmonics
$n = 1,\dots,$ \code{nHarmonics} and the orders $2k = 2, 4, \dots,$ \code{maxOrder}, with and without a gap in
(pseudo)rapidity, and by the scalar product method. Two sets of particles enter: the reference particles (RF), selected
by the particle filters of the task (\code{PARTICLE\_FILTER}, one histogram folder per event filter and RF filter), and
the particles of interest (POI), selected by a second list of filters (\code{POI\_FILTER}), whose differential flow is
measured in every folder versus the event variable and the transverse momentum of the POI.

\paragraph{Correlators.}
With the weighted flow vectors of the RF of an event, $Q_{h,p} = \sum_i w_i^p e^{ih\varphi_i}$, the numerator of the
$m$-particle correlator of harmonics $h_1,\dots,h_m$, a sum over the $m$-tuples of distinct particles, is obtained without
a loop over tuples by summing over the partitions of the $m$ slots into blocks $B$ (the slots given to the same
particle)~\cite{Bilandzic:2013kga}:
\begin{equation}
  N(S) = \sum_{B \ni \min S} (-1)^{|B|-1}(|B|-1)!\; Q_{h_B,|B|}\; N(S\setminus B),
\end{equation}
with $N(\emptyset) = 1$, where $S$ is a set of slots, $h_B$ is the sum of the harmonics of $B$, and $\min S$ is the first slot of $S$.
\code{FlowCorrelator} computes $N(S)$ for all the $2^m$ subsets at once, in $3^m$ operations. With the harmonics
arranged as $(+n,-n,+n,-n,\dots)$, the first $2k$ slots give the correlator $\avg{2k}_n = N/D$, where the denominator
$D$ is the same sum with all harmonics zero. One call therefore gives all the orders. The event-averaged correlators
$\avg{\avg{2k}}_n$ are profiles weighted by $D$. For the differential correlators $\avg{2k'}_n$, a POI takes the first
slot and RF the others. The block of the first slot uses the vectors $p_{h}$ of the POI, or, when it holds more than one
slot, the vectors $q_{h,p}$ of the particles that are both POI and RF~\cite{Bilandzic:2010jr}:
\begin{equation}
  N'(S) = \sum_{T\subseteq S\setminus\{1\}} (-1)^{|T|}\,|T|!\; X_{h_1+h_T,\,|T|+1}\; N(S\setminus\{1\}\setminus T),
\end{equation}
with $X_{h,1} = p_h$ and $X_{h,p} = q_{h,p}$ for $p\ge 2$. Self-correlations are thus removed exactly, whatever the overlap
of the two selections. The particle weights $w_i$ (the inverse of the efficiencies of Section~\ref{sec:efficiency})
correct for a non-uniform acceptance.

\paragraph{Gap and sub-events.}
With \code{gap} on, the RF are divided into two sub-events, $A$ ($x<-\Delta/2$) and $B$ ($x>\Delta/2$), where $x$ is the
pseudorapidity or the rapidity (\code{gapVariable}) and $\Delta$ is \code{gapSize}. The gapped correlator takes $k$
particles of harmonic $n$ in $A$ and $k$ of harmonic $-n$ in $B$:
$N_A(n,\dots,n)\,N_B(n,\dots,n)^*$~\cite{Jia:2017hbm}. Non-flow correlations within a jet or a resonance decay, which do
not reach across the gap, are suppressed. For the POI, a POI and $k-1$ RF are taken in one sub-event and $k$ RF in the
other, and both assignments are summed.

\paragraph{Cumulants and flow.}
The derived calculator obtains the cumulants $c_n\{2k\}$ from the generating function~\cite{Borghini:2001vi},
\begin{equation}
  \sum_{k\ge1} c_n\{2k\}\frac{z^k}{(k!)^2} = \ln\Bigl(1+\sum_{k\ge1}\avg{\avg{2k}}_n\frac{z^k}{(k!)^2}\Bigr),
\end{equation}
by the recursion of the logarithm of a series, for any order. The first orders are
\begin{align}
  c_n\{2\} &= \avg{\avg{2}}, \nonumber\\
  c_n\{4\} &= \avg{\avg{4}} - 2\avg{\avg{2}}^2, \nonumber\\
  c_n\{6\} &= \avg{\avg{6}} - 9\avg{\avg{4}}\avg{\avg{2}} + 12\avg{\avg{2}}^3, \\
  c_n\{8\} &= \avg{\avg{8}} - 16\avg{\avg{6}}\avg{\avg{2}} - 18\avg{\avg{4}}^2 \nonumber\\
           &\quad + 144\avg{\avg{4}}\avg{\avg{2}}^2 - 144\avg{\avg{2}}^4. \nonumber
\end{align}
For a single flow magnitude $v$, $c_n\{2k\} = A_k v^{2k}$ with $A_k = 1, -1, 4, -33, 456, \dots$, so that
$v_n\{2k\} = (c_n\{2k\}/A_k)^{1/2k}$, which is defined when $c_n\{2k\}/A_k > 0$. The differential cumulant $d_n\{2k\}$ is
the cumulant with the block of the POI replaced, that is, the derivative of $c_n\{2k\}$ along
$\avg{\avg{2j}}\to\avg{\avg{2j}}+(j/k)\avg{\avg{2j'}}$:
\begin{align}
  d_n\{4\} &= \avg{\avg{4'}} - 2\avg{\avg{2'}}\avg{\avg{2}}, \nonumber\\
  d_n\{6\} &= \avg{\avg{6'}} - 6\avg{\avg{4'}}\avg{\avg{2}} - 3\avg{\avg{2'}}\avg{\avg{4}} \\
           &\quad + 12\avg{\avg{2'}}\avg{\avg{2}}^2, \nonumber
\end{align}
and $v'_n\{2k\} = d_n\{2k\}/(A_k\, v_n\{2k\}^{2k-1})$. The scalar product method~\cite{Adler:2002pu} uses the sub-events
of the gap: $v_n\{\mathrm{SP}\} = \avg{u_n Q^*}/\sqrt{\avg{Q_A Q_B^*}}$, where $u_n = e^{in\varphi}$ is the unit vector of
a POI and $Q$ is the flow vector of the opposite sub-event. The uncertainties of the derived histograms are propagated
linearly from those of the correlators, without their covariances; the sub-samples (Section~\ref{sec:subsample}) give
the full uncertainties. Table~\ref{tab:flow} lists the histograms. \code{BasicEventGen} can impose flow coefficients on
its particles (\code{Flow:Vn}, constant or rising with $\pT$), which serve to test the analysis. A test of the module
recovers $v_2 = 0.1$ and $v_3 = 0.05$ at all orders up to 8, with a hole in the acceptance corrected by weights.

\begin{table*}[t]
  \centering
  \caption{Histograms of a flow folder (\code{ParticleFlowHistos}), for the harmonic $n$, the order $2k$ and the POI
  filter $j$, as functions of the event variable; the names ending in \code{\_pt} are functions of the event variable
  and of the $\pT$ of the POI.}
  \label{tab:flow}
  \small
  \begin{tabular}{@{}lll@{}}
    \toprule
    Name & Content & Option \\
    \midrule
    \multicolumn{3}{@{}l}{\emph{Base (profiles)}} \\
    \code{evt}                                      & distribution of the event variable & \\
    \code{Corr}$2k$\code{\_n}$n$                    & $\avg{\avg{2k}}_n$ of the RF & \\
    \code{CorrGap}$2k$\code{\_n}$n$                 & $\avg{\avg{2k}}_n$ across the gap & \code{gap} \\
    \code{SPQaQb\_n}$n$                             & $\avg{Q_AQ_B^*}_n$ & \code{scalarProduct} \\
    \code{POI}$j$\code{\_Corr}$2k$\code{\_n}$n$[\code{\_pt}]    & $\avg{\avg{2k'}}_n$ & \\
    \code{POI}$j$\code{\_CorrGap}$2k$\code{\_n}$n$[\code{\_pt}] & $\avg{\avg{2k'}}_n$ across the gap & \code{gap} \\
    \code{POI}$j$\code{\_SPuQ\_n}$n$[\code{\_pt}]   & $\avg{u_nQ^*}$ & \code{scalarProduct} \\
    \midrule
    \multicolumn{3}{@{}l}{\emph{Derived}} \\
    \code{Cum}$2k$\code{\_n}$n$, \code{V}$2k$\code{\_n}$n$       & $c_n\{2k\}$, $v_n\{2k\}$ & \\
    \code{CumGap}$2k$\code{\_n}$n$, \code{VGap}$2k$\code{\_n}$n$ & the same across the gap & \code{gap} \\
    \code{POI}$j$\code{\_DCum}$2k$\code{\_n}$n$[\code{\_pt}], \code{POI}$j$\code{\_V}$2k$\code{\_n}$n$[\code{\_pt}] & $d_n\{2k\}$, $v'_n\{2k\}$ & \\
    \code{POI}$j$\code{\_DCumGap}$2k$\ldots, \code{POI}$j$\code{\_VGap}$2k$\ldots & the same across the gap & \code{gap} \\
    \code{POI}$j$\code{\_VSP\_n}$n$[\code{\_pt}]    & $v_n\{\mathrm{SP}\}$ & \code{scalarProduct} \\
    \bottomrule
  \end{tabular}
\end{table*}

% ---------------------------------------------------------------------------------------------------------------
\subsection{Jets}
\label{sec:jets}

The module Jets, built when \CAP{} is configured with FastJet~\cite{Cacciari:2011ma}, finds the jets of the events and
measures their properties. The task \code{JetAnalyzer} clusters the particles accepted by each of its particle filters
with the anti-$k_T$~\cite{Cacciari:2008gp}, $k_T$ or Cambridge/Aachen algorithm, for each radius $R$ of a list given in
one run (\code{JET:R}, e.g. \code{0.2, 0.4, 0.6}). The jets kept are those above a minimum transverse momentum
(\code{JET:PT\_MIN}) and fully inside the acceptance of the particles, $|\eta_{\rm jet}| < \eta_{\max} - R$
(\code{JET:PARTICLE\_ETA\_MAX}). They are ranked by decreasing $p_T$, the leading jet having rank 1. The jets of an event
are found once per particle filter and radius, and kept for the event, so that every jet task of the job uses the same
jets without clustering them again (\code{JetClustering}).

Jet filters (\code{JET\_FILTER}) select the jets, with conditions on their transverse momentum, pseudorapidity,
azimuth, mass, number of constituents, and rank in the event: the leading jet, the subleading jet, or any range of
ranks. They are defined in the library of RunCAP, as the event and particle filters, and are created by the task
\code{JetFilterCreator}. The histograms of \code{JetAnalyzer} fill one folder per event filter, particle filter,
radius and jet filter, named for instance \code{MB\_Charged\_R0p4\_Leading}. They hold:
\begin{itemize}
  \item per event: the number of jets accepted and the scalar sum of their $p_T$ (events without a jet included);
  \item per jet: $p_T$, $\eta$, $\varphi$, the mass, the numbers of constituents in total and by charge sign
  ($N$, $N_+$, $N_-$, $N_0$), the net charge, and the jet charge $Q_\kappa = \sum_i q_i\,(p_{T,i}/p_{T,{\rm jet}})^\kappa$;
  their averages are also profiles versus the jet $p_T$;
  \item per constituent, in total and by charge sign: the momentum fraction $z = p_{T,i}/p_{T,{\rm jet}}$, the
  momentum $j_T$ perpendicular to the jet axis, the angle $\theta$ to the axis, the distance $r$ to the axis in rapidity
  and azimuth, and the transverse momentum profile $\rho(r)$.
\end{itemize}
The per-event histograms are normalised by the number of events accepted by the event filter. The per-jet and
per-constituent histograms characterise the jets themselves: they are normalised by the number of jets of their folder,
that is of their event filter, particle filter, radius and jet filter, so that a constituent distribution is a
distribution per jet. The profiles are averages per jet.
Two further tasks will use the same jets: an analysis of jet pairs (filters on the ranks of the jets, their
azimuthal separation, acoplanarity, momentum balance and mass), and an analysis of the fluctuations and correlations
of the multiplicities within the jets ($\avg{N}$, variances and second factorial moments by charge, $R_2$, $\nu_{\rm dyn}$),
corrected for the underlying event measured in cones transverse to the leading jet.

% ---------------------------------------------------------------------------------------------------------------
\subsection{Global observables and event shapes}
\label{sec:global}

The task \code{GlobalAnalyzer} computes, for each event and each particle filter, the multiplicity $n$, the total energy
$e=\sum_i E_i$, the net electric charge $q=\sum_i q_i$, the net strangeness $s$ and the net baryon number $b$ of the
accepted particles, their total transverse momentum $\sum\pT$ and its mean $\sum\pT/n$. They are the global quantities
of the selections of its particle filters, made once per event and shared by the tasks (Section~\ref{sec:filters}), which
the event filters of any task also test directly to define event classes (multiplicity classes, for instance);
\code{NuDynAnalyzer} uses one of them as its event variable. The histograms hold their event-by-event
distributions and their averages versus $n$ (Table~\ref{tab:global}); the derived calculator adds the inclusive mean
transverse momentum versus multiplicity, $\avg{\sum\pT}/n$.

The task \code{SpherocityAnalyzer} computes the transverse spherocity of the accepted particles of each event,
\begin{equation}
  S_0 = \frac{\pi^2}{4}\,\min_{\hat{\mathbf n}} \left(\frac{\sum_i |\mathbf p_{\mathrm{T},i}\times\hat{\mathbf n}|}{\sum_i p_{\mathrm{T},i}}\right)^2 ,
\end{equation}
where $\hat{\mathbf n}$ is a unit vector of the transverse plane, and its unweighted version $S_1$, computed with the
unit vectors $\mathbf p_{\mathrm{T},i}/p_{\mathrm{T},i}$. Both range from 0 for pencil-like events to 1 for isotropic
ones. Since the sum is a concave function of the azimuth of $\hat{\mathbf n}$ between the directions of the particles,
its minimum lies along one of them: these directions are tried, which gives the exact minimum, for up to \code{nSteps}
particles, and \code{nSteps} directions spread over $\pi$ are tried for more. Selections of fewer than
\code{MinMultiplicity} particles (by default 3) are not filled, since their spherocity is not meaningful; analyses of
data usually require at least 10 charged particles. The same quantity is that of the selections tested by the event
filters (Section~\ref{sec:filters}).
\todo{reference for spherocity.}

\begin{table*}[t]
  \centering
  \caption{Histograms of the folders of global observables (\code{GlobalHistos}) and of spherocity
  (\code{SpherocityHistos}). Distributions and profiles are not scaled.}
  \label{tab:global}
  \small
  \begin{tabular}{@{}lll@{}}
    \toprule
    Name & Content & Option \\
    \midrule
    \multicolumn{3}{@{}l}{\emph{Global, base}} \\
    \code{n}, \code{e}, \code{q}, \code{s}, \code{b} & distributions of $n$, $e$, $q$, $s$, $b$ & \\
    \code{ptSum}, \code{ptAvg}   & distributions of $\sum\pT$, $\sum\pT/n$ & \\
    \code{n\_}$v$\code{Prof}       & $\avg{v}$ vs $n$, $v$ = \code{e}, \code{q}, \code{s}, \code{b}, \code{ptSum}, \code{ptAvg} & \\
    \multicolumn{3}{@{}l}{\emph{Global, derived}} \\
    \code{n\_ptAvgIncl}          & $\avg{\sum\pT}/n$ vs $n$ & \\
    \midrule
    \multicolumn{3}{@{}l}{\emph{Spherocity}} \\
    \code{S0}, \code{S1}         & distributions of $S_0$, $S_1$ & \code{FillS0}, \code{FillS1} \\
    \code{S1VsS0}                & $S_1$ vs $S_0$                & \code{FillS1VsS0} \\
    \bottomrule
  \end{tabular}
\end{table*}

% ---------------------------------------------------------------------------------------------------------------
\subsection{Performance analysis}
\label{sec:performance}

For performance studies, \code{MeasurementSimulator} produces from each generated event a detector-level event, with
the efficiencies and smearers of Section~\ref{sec:efficiency}, and \code{PerformanceAnalyzer} compares the two: for each
particle, the differences between its smeared and generated kinematics, $\Delta\pT$, $\Delta\eta$, $\Delta y$ and
$\Delta\varphi$, are histogrammed, and their averages are recorded as profiles versus $\pT$ and $\eta$
(Table~\ref{tab:performance}). The same analyses, run on the generated and on the detector-level events, give the
generator-level and the measured observables; \code{ClosureCalculator} takes their differences or ratios, histogram by
histogram, to assess the effect of the detector on each observable.

\begin{table*}[t]
  \centering
  \caption{Histograms of a performance folder (\code{ParticlePerformanceHistos}); $\Delta$ is detector level minus
  generator level.}
  \label{tab:performance}
  \small
  \begin{tabular}{@{}ll@{}}
    \toprule
    Name & Content \\
    \midrule
    \code{n1\_dpt}, \code{n1\_deta}, \code{n1\_dy}, \code{n1\_dphi} & distributions of $\Delta\pT$, $\Delta\eta$, $\Delta y$, $\Delta\varphi$ \\
    \code{d}$v$\code{VsPt}   & $\avg{\Delta v}$ vs $\pT$, $v$ = \code{pt}, \code{eta}, \code{y}, \code{phi} \\
    \code{d}$v$\code{VsEta}  & $\avg{\Delta v}$ vs $\eta$ \\
    \bottomrule
  \end{tabular}
\end{table*}
